%--------------------------------------------------------------------%
%
% Halaman Judul lembar kedua (Pedoman UM Lampiran 3a)
%
%--------------------------------------------------------------------%

\clearpage
\thispagestyle{empty}

\begin{center}
    \setstretch{1.5}

    \bfseries \MakeUppercase{\judul}
    \vfill

    \normalfont
    SKRIPSI \\[0.3\baselineskip]
    Diajukan kepada \\
    Universitas Negeri Malang \\
    untuk memenuhi salah satu persyaratan \\
    dalam menyelesaikan program \gelarprogram
    \vfill

    OLEH \\
    \MakeUppercase{\penulis} \\
    NIM \nim
    \vfill

    % ponytail: placeholder logo Ganesha. GANTI dengan logo UM (diameter 3 cm).
    \includegraphics[width=3cm]{resources/cover-ganesha.jpg}
    \vfill

    \bfseries
    \MakeUppercase{Universitas Negeri Malang} \\
    \MakeUppercase{\fakultas} \\
    \MakeUppercase{\programstudi} \\
    \MakeUppercase{\bulantahun}

\end{center}

\clearpage
